\section{Ringed spaces}
\label{section:0.4}

\subsection{Ringed spaces, $\mathcal{A}$-modules, $\mathcal{A}$-algebras}
\label{subsection:0.4.1}

\begin{env}[4.1.1]
\label{0.4.1.1}
A \emph{ringed space} (resp. \emph{topologically ringed space}) is a pair
$(X,\mathcal{A})$ consisting of a topological space $X$ and a sheaf of rings
(not necessarily commutative) (resp. a sheaf of topological rings)
$\mathcal{A}$ on $X$; one says that $X$ is the
\oldpage[0\textsubscript{I}]{36}
topological space \emph{underlying} the ringed space $(X,\mathcal{A})$, and
$\mathcal{A}$ its \emph{structure sheaf}. The latter is denoted by
$\mathcal{O}_X$, and its stalk at a point $x\in X$ is denoted by $\mathcal{O}_{X,x}$,
or simply $\mathcal{O}_x$ when no confusion results.

We shall denote by $1$ or $e$ the \emph{unit section} of $\mathcal{O}_X$
over $X$ (the unit element of $\Gamma(X,\mathcal{O}_X)$).

Since in this Treatise we shall mostly have to consider sheaves of
\emph{commutative} rings, it will be understood, when we speak of a ringed
space $(X,\mathcal{A})$ without further specification, that $\mathcal{A}$ is a sheaf
of commutative rings.

Ringed spaces whose structure sheaf is not necessarily commutative
(resp. topologically ringed spaces) form a \emph{category}
when one defines a \emph{morphism}
$(X,\mathcal{A})\to(Y,\mathcal{B})$ as a pair
$(\psi,\theta)=\Psi$ consisting of a continuous map $\psi:X\to Y$ and a
$\psi$-\emph{morphism} $\theta:\mathcal{B}\to\mathcal{A}$ (3.5.1) of
sheaves of rings (resp. of sheaves of topological rings); the
\emph{composite} of a second morphism
$\Psi'=(\psi',\theta'):(Y,\mathcal{B})\to(Z,\mathcal{C})$ and $\Psi$, denoted by
$\Psi''=\Psi'\circ\Psi$, is the morphism $(\psi'',\theta'')$, where
$\psi''=\psi'\circ\psi$ and $\theta''$ is the composite of $\theta$ and $\theta'$
(equal to $\psi'_*(\theta)\circ\theta'$, cf. 3.5.2). For ringed
spaces, recall that one then has
$\theta''^\sharp=\theta^\sharp\circ\psi^*(\theta'^\sharp)$ (3.5.5); hence, if
$\theta'^\sharp$ and $\theta^\sharp$ are \emph{injective}
(resp. \emph{surjective}) homomorphisms, so is $\theta''^\sharp$, in view
of the fact that $\psi_x\circ\rho_{\psi(x)}$ is an isomorphism for every $x\in X$
(3.7.2). Using the foregoing, one verifies at once that when $\psi$
is an \emph{injective} continuous map and $\theta^\sharp$ a
\emph{surjective} homomorphism of sheaves of rings, the morphism
$(\psi,\theta)$ is a \emph{monomorphism} (T, 1.1) in the category of
ringed spaces.

By abuse of language, we shall often replace $\psi$ by $\Psi$ in the
notation, for example writing $\Psi^{-1}(U)$ instead of $\psi^{-1}(U)$
for a subset $U$ of $Y$, when this is not likely to cause confusion.
\end{env}

\begin{env}[4.1.2]
\label{0.4.1.2}
For every subset $M$ of $X$, the pair $(M,\mathcal{A}|M)$ is evidently a
ringed space, said to be \emph{induced} on $M$ by the ringed space
$(X,\mathcal{A})$ (and also called the \emph{restriction} of
$(X,\mathcal{A})$ to $M$). If $j$ is the canonical injection $M\to X$ and
$\omega$ the identity map of $\mathcal{A}|M$, then $(j,\omega^b)$ is a
monomorphism $(M,\mathcal{A}|M)\to(X,\mathcal{A})$ of ringed spaces, called
the \emph{canonical injection}. The composite of a morphism
$\Psi:(X,\mathcal{A})\to(Y,\mathcal{B})$ and this injection is called the
\emph{restriction} of $\Psi$ to $M$.
\end{env}

\begin{env}[4.1.3]
\label{0.4.1.3}
We shall not return to the definition of $\mathcal{A}$-modules, or
algebraic sheaves, on a ringed space $(X,\mathcal{A})$
(G, II, 2.2); when $\mathcal{A}$ is a sheaf of not necessarily
commutative rings, by $\mathcal{A}$-module one must always
understand ``left $\mathcal{A}$-module'' unless the contrary is
expressly stated. The sub-$\mathcal{A}$-modules of $\mathcal{A}$ will be called
sheaves of (left, right or two-sided) ideals in $\mathcal{A}$, or
$\mathcal{A}$-Ideals.

When $\mathcal{A}$ is a sheaf of commutative rings and, in the
definition of $\mathcal{A}$-modules, one replaces the module structure
everywhere by that of algebra, one obtains the definition of an
$\mathcal{A}$-algebra on $X$. It amounts to the same to say that a
(not necessarily commutative) $\mathcal{A}$-algebra is an
$\mathcal{A}$-module $\mathcal{C}$ endowed with a homomorphism of
$\mathcal{A}$-modules
$\varphi:\mathcal{C}\otimes_{\mathcal{A}}\mathcal{C}\to\mathcal{C}$ and
a section $e$ over $X$ such that: 1\textsuperscript{o} the diagram
\oldpage[0\textsubscript{I}]{37}
\[
\begin{tikzcd}[column sep=large,row sep=large]
\mathcal{C}\otimes_{\mathcal{A}}\mathcal{C}\otimes_{\mathcal{A}}\mathcal{C}
  \arrow[r,"\varphi\otimes 1"]
  \arrow[d,"1\otimes\varphi"']
&
\mathcal{C}\otimes_{\mathcal{A}}\mathcal{C}
  \arrow[d,"\varphi"]
\\
\mathcal{C}\otimes_{\mathcal{A}}\mathcal{C}
  \arrow[r,"\varphi"']
&
\mathcal{C}
\end{tikzcd}
\]
is commutative; 2\textsuperscript{o} for every open set $U\subset X$ and every section
$s\in\Gamma(U,\mathcal{C})$, one has
$\varphi((e|U)\otimes s)=\varphi(s\otimes(e|U))=s$. To say that
$\mathcal{C}$ is a commutative $\mathcal{A}$-algebra amounts, moreover, to
saying that the diagram
\[
\begin{tikzcd}[column sep=large,row sep=large]
\mathcal{C}\otimes_{\mathcal{A}}\mathcal{C}
  \arrow[rr,"\sigma"]
  \arrow[dr,"\varphi"']
&&
\mathcal{C}\otimes_{\mathcal{A}}\mathcal{C}
  \arrow[dl,"\varphi"]
\\
& \mathcal{C} &
\end{tikzcd}
\]
is commutative, where $\sigma$ denotes the canonical symmetry of the tensor product
$\mathcal{C}\otimes_{\mathcal{A}}\mathcal{C}$.

Homomorphisms of $\mathcal{A}$-algebras are also defined like the
homomorphisms of $\mathcal{A}$-modules in (G, II, 2.2), but naturally
they no longer form an abelian group.

If $\mathcal{M}$ is a sub-$\mathcal{A}$-module of an
$\mathcal{A}$-algebra $\mathcal{C}$, the sub-$\mathcal{A}$-algebra of
$\mathcal{C}$ generated by $\mathcal{M}$ is the sum of the images of the
homomorphisms $\bigotimes^n\mathcal{M}\to\mathcal{C}$ (for
$n\geqslant 0$). It is also the sheaf associated to the presheaf of algebras
$U\mapsto\mathcal{B}(U)$, where $\mathcal{B}(U)$ is the subalgebra
of $\Gamma(U,\mathcal{C})$ generated by the submodule
$\Gamma(U,\mathcal{M})$.
\end{env}

\begin{env}[4.1.4]
\label{0.4.1.4}
A sheaf of rings $\mathcal{A}$ on a topological space $X$ is said to be
\emph{reduced at a point} $x$ of $X$ if the stalk $\mathcal{A}_x$ is a
\emph{reduced} ring (1.1.1); $\mathcal{A}$ is said to be \emph{reduced}
if it is reduced at every point of $X$. Recall that a ring $A$ is said to be
\emph{regular} if each of the local rings $A_{\mathfrak p}$ (where
$\mathfrak p$ ranges over the set of prime ideals of $A$) is a regular
local ring; we shall say that a sheaf of rings $\mathcal{A}$ on $X$
is \emph{regular at a point} $x$ (resp. \emph{regular}) if the stalk
$\mathcal{A}_x$ is a regular ring (resp. if $\mathcal{A}$ is regular
at every point). Finally, we shall say that a sheaf of rings $\mathcal{A}$ on
$X$ is \emph{normal at a point} $x$ (resp. \emph{normal}) if its stalk
$\mathcal{A}_x$ is an integral and integrally closed ring (resp. if
$\mathcal{A}$ is normal at every point). We shall say that a ringed space
$(X,\mathcal{A})$ has one of the preceding properties if the sheaf of rings
$\mathcal{A}$ has that property.

A sheaf of \emph{graded} rings $\mathcal{A}$ is by definition a
sheaf of rings that is the direct sum (G, II, 2.7) of a family
$(\mathcal{A}_n)_{n\in\mathbf Z}$ of sheaves of abelian groups satisfying the
conditions $\mathcal{A}_m\mathcal{A}_n\subset\mathcal{A}_{m+n}$; a
\emph{graded $\mathcal{A}$-module} is an $\mathcal{A}$-module
$\mathcal{F}$ that is the direct sum of a family
$(\mathcal{F}_n)_{n\in\mathbf Z}$ of sheaves of abelian groups,
satisfying the conditions
$\mathcal{A}_m\mathcal{F}_n\subset\mathcal{F}_{m+n}$. Moreover, it amounts
to the same to say that
$(\mathcal{A}_m)_x(\mathcal{A}_n)_x\subset(\mathcal{A}_{m+n})_x$
(resp.
$(\mathcal{A}_m)_x(\mathcal{F}_n)_x\subset(\mathcal{F}_{m+n})_x$)
at every point $x$.
\end{env}

\begin{env}[4.1.5]
\label{0.4.1.5}
Given a (not necessarily commutative) ringed space $(X,\mathcal{A})$,
we shall not recall here the definitions of the bifunctors
$\mathcal{F}\otimes_{\mathcal{A}}\mathcal{G}$,
$\mathcal{H}\!om_{\mathcal{A}}(\mathcal{F},\mathcal{G})$ and
$\operatorname{Hom}_{\mathcal{A}}(\mathcal{F},\mathcal{G})$
(G, II, 2.8 and 2.2) on the categories of left or
right $\mathcal{A}$-modules (as the case may be), with values in the category of sheaves
of abelian groups (or more

\oldpage[0\textsubscript{I}]{38}
generally of $\mathcal{C}$-modules, if $\mathcal{C}$ is the centre of
$\mathcal{A}$). The stalk
$(\mathcal{F}\otimes_{\mathcal{A}}\mathcal{G})_x$ at every point $x\in X$
is canonically identified with
$\mathcal{F}_x\otimes_{\mathcal{A}_x}\mathcal{G}_x$, and one defines a
canonical functorial homomorphism
\[
  (\mathcal{H}\!om_{\mathcal{A}}(\mathcal{F},\mathcal{G}))_x
  \longrightarrow
  \operatorname{Hom}_{\mathcal{A}_x}(\mathcal{F}_x,\mathcal{G}_x)
\]
which, in general, is neither injective nor surjective. The bifunctors considered
above are additive and, in particular, commute with finite direct
sums; $\mathcal{F}\otimes_{\mathcal{A}}\mathcal{G}$ is right exact in
$\mathcal{F}$ and in $\mathcal{G}$, commutes with inductive limits, and
$\mathcal{A}\otimes_{\mathcal{A}}\mathcal{G}$
(resp. $\mathcal{F}\otimes_{\mathcal{A}}\mathcal{A}$) is
canonically identified with $\mathcal{G}$ (resp. $\mathcal{F}$). The functors
$\mathcal{H}\!om_{\mathcal{A}}(\mathcal{F},\mathcal{G})$ and
$\operatorname{Hom}_{\mathcal{A}}(\mathcal{F},\mathcal{G})$ are
\emph{left exact} in $\mathcal{F}$ and $\mathcal{G}$; more precisely,
if there is an exact sequence of the form
$0\to\mathcal{G}'\to\mathcal{G}\to\mathcal{G}''$, the sequence
\[
  0\longrightarrow
  \mathcal{H}\!om_{\mathcal{A}}(\mathcal{F},\mathcal{G}')
  \longrightarrow
  \mathcal{H}\!om_{\mathcal{A}}(\mathcal{F},\mathcal{G})
  \longrightarrow
  \mathcal{H}\!om_{\mathcal{A}}(\mathcal{F},\mathcal{G}'')
\]
is exact, and if there is an exact sequence of the form
$\mathcal{F}'\to\mathcal{F}\to\mathcal{F}''\to0$, the sequence
\[
  0\longrightarrow
  \mathcal{H}\!om_{\mathcal{A}}(\mathcal{F}'',\mathcal{G})
  \longrightarrow
  \mathcal{H}\!om_{\mathcal{A}}(\mathcal{F},\mathcal{G})
  \longrightarrow
  \mathcal{H}\!om_{\mathcal{A}}(\mathcal{F}',\mathcal{G})
\]
is exact, with analogous properties for the functor
$\operatorname{Hom}$. Moreover,
$\mathcal{H}\!om_{\mathcal{A}}(\mathcal{A},\mathcal{G})$ is
canonically identified with $\mathcal{G}$; finally, for every open set $U\subset X$, one has
\[
  \Gamma\bigl(U,\mathcal{H}\!om_{\mathcal{A}}(\mathcal{F},\mathcal{G})\bigr)
  =
  \operatorname{Hom}_{\mathcal{A}|U}(\mathcal{F}|U,\mathcal{G}|U).
\]

For every left (resp. right) $\mathcal{A}$-module $\mathcal{F}$, the
right (resp. left) $\mathcal{A}$-module
$\mathcal{H}\!om_{\mathcal{A}}(\mathcal{F},\mathcal{A})$ is called the
\emph{dual of $\mathcal{F}$} and is denoted by $\check{\mathcal{F}}$.

Finally, if $\mathcal{A}$ is a sheaf of commutative rings and
$\mathcal{F}$ an $\mathcal{A}$-module, then
$U\mapsto\bigwedge^p\Gamma(U,\mathcal{F})$ is a presheaf whose
associated sheaf is an $\mathcal{A}$-module, which is denoted by
$\bigwedge^p\mathcal{F}$ and called the \emph{$p$-th exterior power
of $\mathcal{F}$}; one easily verifies that the canonical map from the
presheaf $U\mapsto\bigwedge^p\Gamma(U,\mathcal{F})$ to the
associated sheaf $\bigwedge^p\mathcal{F}$ is \emph{injective}, and that for every
$x\in X$ one has
$(\bigwedge^p\mathcal{F})_x=\bigwedge^p(\mathcal{F}_x)$. It is clear that
$\bigwedge^p\mathcal{F}$ is a covariant functor in $\mathcal{F}$.
\end{env}

\begin{env}[4.1.6]
\label{0.4.1.6}
Suppose that $\mathcal{A}$ is a sheaf of not necessarily
commutative rings, $\mathcal{J}$ a sheaf of left ideals of $\mathcal{A}$,
and $\mathcal{F}$ a left $\mathcal{A}$-module; one then denotes by
$\mathcal{J}\mathcal{F}$ the sub-$\mathcal{A}$-module of $\mathcal{F}$
that is the image of $\mathcal{J}\otimes_{\mathbf Z}\mathcal{F}$ (where $\mathbf Z$ is the
sheaf associated to the constant presheaf $U\to\mathbf Z$) under the
canonical map $\mathcal{J}\otimes_{\mathbf Z}\mathcal{F}\to\mathcal{F}$; it is
clear that for every $x\in X$ one has
$(\mathcal{J}\mathcal{F})_x=\mathcal{J}_x\mathcal{F}_x$. When
$\mathcal{A}$ is commutative, $\mathcal{J}\mathcal{F}$ is also the image
of the canonical map $\mathcal{J}\otimes_{\mathcal{A}}\mathcal{F}\to\mathcal{F}$.
It is immediate that $\mathcal{J}\mathcal{F}$ is also the
$\mathcal{A}$-module associated to the presheaf
$U\to\Gamma(U,\mathcal{J})\Gamma(U,\mathcal{F})$. If
$\mathcal{J}_1$, $\mathcal{J}_2$ are two sheaves of left ideals of
$\mathcal{A}$, one has
$\mathcal{J}_1(\mathcal{J}_2\mathcal{F})=(\mathcal{J}_1\mathcal{J}_2)\mathcal{F}$.
\end{env}

\begin{env}[4.1.7]
\label{0.4.1.7}
Let $(X_\lambda,\mathcal{A}_\lambda)_{\lambda\in L}$ be a family of ringed
spaces; for every pair $(\lambda,\mu)$, suppose given an open subset
$V_{\lambda\mu}$ of $X_\lambda$ and an isomorphism of ringed spaces
\[
  \varphi_{\lambda\mu}:
  (V_{\mu\lambda},\mathcal{A}_\mu|V_{\mu\lambda})
  \xrightarrow{\sim}
  (V_{\lambda\mu},\mathcal{A}_\lambda|V_{\lambda\mu}),
\]
with $V_{\lambda\lambda}=X_\lambda$, $\varphi_{\lambda\lambda}$ being
the identity. Suppose moreover that, for every triple $(\lambda,\mu,\nu)$, if
$\varphi'_{\mu\lambda}$ denotes the restriction of
$\varphi_{\mu\lambda}$ to $V_{\lambda\mu}\cap V_{\lambda\nu}$, then
$\varphi'_{\mu\lambda}$ is an isomorphism of
\[
  (V_{\lambda\mu}\cap V_{\lambda\nu},
   \mathcal{A}_\lambda|(V_{\lambda\mu}\cap V_{\lambda\nu}))
\]
onto
\[
  (V_{\mu\nu}\cap V_{\mu\lambda},
   \mathcal{A}_\mu|(V_{\mu\nu}\cap V_{\mu\lambda}))
\]
and that
$\varphi'_{\lambda\nu}=\varphi'_{\lambda\mu}\circ\varphi'_{\mu\nu}$
(\emph{gluing condition} for the $\varphi_{\lambda\mu}$). One can then
first consider the topological space obtained by gluing (by
means of the $\varphi_{\lambda\mu}$) the $X_\lambda$

\oldpage[0\textsubscript{I}]{39}
along the $V_{\lambda\mu}$; if $X_\lambda$ is identified with the corresponding open
subset $X'_\lambda$ of $X$, the hypotheses imply that the three
sets $V_{\lambda\mu}\cap V_{\lambda\nu}$,
$V_{\mu\nu}\cap V_{\mu\lambda}$,
$V_{\nu\lambda}\cap V_{\nu\mu}$ are identified with
$X'_\lambda\cap X'_\mu\cap X'_\nu$. One can then transport to
$X'_\lambda$ the ringed space structure of $X_\lambda$, and if
$\mathcal{A}'_\lambda$ is the sheaf of rings transported from
$\mathcal{A}_\lambda$, the $\mathcal{A}'_\lambda$ satisfy the gluing
condition (3.3.1) and therefore define a sheaf of rings $\mathcal{A}$ on
$X$; $(X,\mathcal{A})$ is called the ringed space obtained by
\emph{gluing the $(X_\lambda,\mathcal{A}_\lambda)$ along the
$V_{\lambda\mu}$}, by means of the $\varphi_{\lambda\mu}$.
\end{env}

\subsection{Direct image of an $\mathcal{A}$-module}
\label{subsection:0.4.2}

\begin{env}[4.2.1]
\label{0.4.2.1}
Let $(X,\mathcal{A})$, $(Y,\mathcal{B})$ be two ringed spaces, and
$\Psi=(\psi,\theta)$ a morphism
$(X,\mathcal{A})\to(Y,\mathcal{B})$; thus $\psi_*(\mathcal{A})$ is a
sheaf of rings on $Y$, and $\theta$ a homomorphism
$\mathcal{B}\to\psi_*(\mathcal{A})$ of sheaves of rings. Now let
$\mathcal{F}$ be an $\mathcal{A}$-module; its direct image
$\psi_*(\mathcal{F})$ is a sheaf of abelian groups on $Y$. Moreover,
for every open set $U\subset Y$,
\[
  \Gamma(U,\psi_*(\mathcal{F}))
  =\Gamma(\psi^{-1}(U),\mathcal{F})
\]
is endowed with a module structure over the ring
$\Gamma(U,\psi_*(\mathcal{A}))=\Gamma(\psi^{-1}(U),\mathcal{A})$; the
bilinear maps defining these structures, being compatible
with the restriction operations, define on $\psi_*(\mathcal{F})$ a
structure of $\psi_*(\mathcal{A})$-module. The homomorphism
$\theta:\mathcal{B}\to\psi_*(\mathcal{A})$ then also makes it possible to define on
$\psi_*(\mathcal{F})$ a structure of $\mathcal{B}$-module; we shall say that
this $\mathcal{B}$-module is the \emph{direct image of $\mathcal{F}$ under the
morphism $\Psi$}, and we shall denote it by $\Psi_*(\mathcal{F})$. If
$\mathcal{F}_1$, $\mathcal{F}_2$ are two $\mathcal{A}$-modules on $X$ and
$u$ is an $\mathcal{A}$-homomorphism $\mathcal{F}_1\to\mathcal{F}_2$, it is
immediate (by considering the sections over the open sets of $Y$) that
$\psi_*(u)$ is a $\psi_*(\mathcal{A})$-homomorphism
$\psi_*(\mathcal{F}_1)\to\psi_*(\mathcal{F}_2)$, and \emph{a fortiori} a
$\mathcal{B}$-homomorphism
$\Psi_*(\mathcal{F}_1)\to\Psi_*(\mathcal{F}_2)$; as a
$\mathcal{B}$-homomorphism, we shall denote it by $\Psi_*(u)$. We thus see that
$\Psi_*$ is a \emph{covariant functor} from the category of
$\mathcal{A}$-modules to that of $\mathcal{B}$-modules. Moreover, it is
immediate that this functor is \emph{left exact} (G, II, 2.12).

On $\psi_*(\mathcal{A})$, the $\mathcal{B}$-module structure and the
structure of sheaf of rings define a structure of
$\mathcal{B}$-algebra; this $\mathcal{B}$-algebra will be denoted by
$\Psi_*(\mathcal{A})$.
\end{env}

\begin{env}[4.2.2]
\label{0.4.2.2}
Let $\mathcal{M}$, $\mathcal{N}$ be two $\mathcal{A}$-modules. For every
open set $U$ of $Y$, there is a canonical map
\[
  \Gamma(\psi^{-1}(U),\mathcal{M})
  \times\Gamma(\psi^{-1}(U),\mathcal{N})
  \longrightarrow
  \Gamma(\psi^{-1}(U),\mathcal{M}\otimes_{\mathcal{A}}\mathcal{N})
\]
which is bilinear over the ring
$\Gamma(\psi^{-1}(U),\mathcal{A})=\Gamma(U,\psi_*(\mathcal{A}))$, and
\emph{a fortiori} over $\Gamma(U,\mathcal{B})$; it therefore defines a
homomorphism
\[
  \Gamma(U,\Psi_*(\mathcal{M}))
  \otimes_{\Gamma(U,\mathcal{B})}
  \Gamma(U,\Psi_*(\mathcal{N}))
  \longrightarrow
  \Gamma(U,\Psi_*(\mathcal{M}\otimes_{\mathcal{A}}\mathcal{N}))
\]
and since one verifies at once that these homomorphisms are compatible with
the restriction operations, they give a canonical functorial homomorphism
of $\mathcal{B}$-modules
\[
  \Psi_*(\mathcal{M})\otimes_{\mathcal{B}}\Psi_*(\mathcal{N})
  \longrightarrow
  \Psi_*(\mathcal{M}\otimes_{\mathcal{A}}\mathcal{N}).
  \tag{4.2.2.1}\label{0.4.2.2.1}
\]

\oldpage[0\textsubscript{I}]{40}
which will in general be neither injective nor surjective. If $\mathcal{P}$ is a
third $\mathcal{A}$-module, one verifies at once that the diagram
\[
\begin{tikzcd}[column sep=large,row sep=large]
\Psi_*(\mathcal{M})\otimes_{\mathcal{B}}\Psi_*(\mathcal{N})
  \otimes_{\mathcal{B}}\Psi_*(\mathcal{P})
  \arrow[r]
  \arrow[d]
&
\Psi_*(\mathcal{M}\otimes_{\mathcal{A}}\mathcal{N})
  \otimes_{\mathcal{B}}\Psi_*(\mathcal{P})
  \arrow[d]
\\
\Psi_*(\mathcal{M})\otimes_{\mathcal{B}}
  \Psi_*(\mathcal{N}\otimes_{\mathcal{A}}\mathcal{P})
  \arrow[r]
&
\Psi_*(\mathcal{M}\otimes_{\mathcal{A}}\mathcal{N}
  \otimes_{\mathcal{A}}\mathcal{P})
\end{tikzcd}
\tag{4.2.2.2}\label{0.4.2.2.2}
\]
is commutative.
\end{env}

\begin{env}[4.2.3]
\label{0.4.2.3}
Let $\mathcal{M}$, $\mathcal{N}$ be two $\mathcal{A}$-modules. For every
open set $U\subset Y$, one has by definition
\[
  \Gamma\bigl(\psi^{-1}(U),
    \mathcal{H}\!om_{\mathcal{A}}(\mathcal{M},\mathcal{N})\bigr)
  =
  \operatorname{Hom}_{\mathcal{A}|V}
    (\mathcal{M}|V,\mathcal{N}|V),
\]
where, to lighten the notation, we have set $V=\psi^{-1}(U)$; the map
$u\to\Psi_*(u)$ is a homomorphism
\[
  \operatorname{Hom}_{\mathcal{A}|V}(\mathcal{M}|V,\mathcal{N}|V)
  \longrightarrow
  \operatorname{Hom}_{\mathcal{B}|U}
    (\Psi_*(\mathcal{M})|U,\Psi_*(\mathcal{N})|U)
\]
for the $\Gamma(U,\mathcal{B})$-module structures; these homomorphisms,
being compatible with the restriction operations, therefore define a
canonical functorial homomorphism of $\mathcal{B}$-modules
\[
  \Psi_*\bigl(\mathcal{H}\!om_{\mathcal{A}}(\mathcal{M},\mathcal{N})\bigr)
  \longrightarrow
  \mathcal{H}\!om_{\mathcal{B}}
    (\Psi_*(\mathcal{M}),\Psi_*(\mathcal{N})).
  \tag{4.2.3.1}\label{0.4.2.3.1}
\]
\end{env}

\begin{env}[4.2.4]
\label{0.4.2.4}
If $\mathcal{C}$ is an $\mathcal{A}$-algebra, the composite homomorphism
\[
  \Psi_*(\mathcal{C})\otimes_{\mathcal{B}}\Psi_*(\mathcal{C})
  \longrightarrow
  \Psi_*(\mathcal{C}\otimes_{\mathcal{A}}\mathcal{C})
  \longrightarrow
  \Psi_*(\mathcal{C})
\]
defines on $\Psi_*(\mathcal{C})$ a structure of $\mathcal{B}$-algebra,
as follows from (4.2.2.2). One sees likewise that if $\mathcal{M}$ is a
$\mathcal{C}$-module, $\Psi_*(\mathcal{M})$ is canonically endowed with a
structure of $\Psi_*(\mathcal{C})$-module.
\end{env}

\begin{env}[4.2.5]
\label{0.4.2.5}
Consider in particular the case where $X$ is a \emph{closed} subspace of
$Y$ and $\psi$ is the canonical injection $j:X\to Y$. If
$\mathcal{B}'=\mathcal{B}|X=j^*(\mathcal{B})$ is the restriction of the sheaf
of rings $\mathcal{B}$ to $X$, an $\mathcal{A}$-module $\mathcal{M}$ can
be regarded as a $\mathcal{B}'$-module by means of the homomorphism
$\theta^\sharp:\mathcal{B}'\to\mathcal{A}$; $\Psi_*(\mathcal{M})$ is then
the $\mathcal{B}$-module that induces $\mathcal{M}$ on $X$ and $0$ elsewhere. If
$\mathcal{N}$ is a second $\mathcal{A}$-module,
$\Psi_*(\mathcal{M})\otimes_{\mathcal{B}}\Psi_*(\mathcal{N})$ is therefore identified
with $\Psi_*(\mathcal{M}\otimes_{\mathcal{B}'}\mathcal{N})$, and
$\mathcal{H}\!om_{\mathcal{B}}(\Psi_*(\mathcal{M}),\Psi_*(\mathcal{N}))$
with
$\Psi_*\bigl(\mathcal{H}\!om_{\mathcal{B}'}(\mathcal{M},\mathcal{N})\bigr)$.
\end{env}

\begin{env}[4.2.6]
\label{0.4.2.6}
Let $(Z,\mathcal{C})$ be a third ringed space, and
$\Psi'=(\psi',\theta')$ a morphism
$(Y,\mathcal{B})\to(Z,\mathcal{C})$; if $\Psi''$ is the composite morphism
$\Psi'\circ\Psi$, it is clear that
$\Psi''_*=\Psi'_*\circ\Psi_*$.
\end{env}

\oldpage[0\textsubscript{I}]{41}
\subsection{Inverse image of a $\mathcal{B}$-module}
\label{subsection:0.4.3}

\begin{env}[4.3.1]
\label{0.4.3.1}
The hypotheses and notation being those of (4.2.1), let $\mathcal{G}$ be a
$\mathcal{B}$-module and $\psi^*(\mathcal{G})$ its inverse image (3.7.1),
which is thus a sheaf of abelian groups on $X$. The definition of the
sections of $\psi^*(\mathcal{G})$ and of $\psi^*(\mathcal{B})$ (3.7.1)
shows that $\psi^*(\mathcal{G})$ is canonically endowed with a structure of
$\psi^*(\mathcal{B})$-module. On the other hand, the homomorphism
$\theta^\sharp:\psi^*(\mathcal{B})\to\mathcal{A}$ endows $\mathcal{A}$ with a
structure of $\psi^*(\mathcal{B})$-module, which we shall denote by
$\mathcal{A}_{[\theta]}$ when confusion must be avoided; the
tensor product
$\psi^*(\mathcal{G})\otimes_{\psi^*(\mathcal{B})}\mathcal{A}_{[\theta]}$
is then endowed with a structure of $\mathcal{A}$-module. We shall say that this
$\mathcal{A}$-module is the inverse image of $\mathcal{G}$ under the morphism
$\Psi$, and we shall denote it by $\Psi^*(\mathcal{G})$. If $\mathcal{G}_1$,
$\mathcal{G}_2$ are two $\mathcal{B}$-modules on $Y$ and $v$ is a
$\mathcal{B}$-homomorphism $\mathcal{G}_1\to\mathcal{G}_2$, then
$\psi^*(v)$, as one verifies at once, is a
$\psi^*(\mathcal{B})$-homomorphism from $\psi^*(\mathcal{G}_1)$ to
$\psi^*(\mathcal{G}_2)$; consequently $\psi^*(v)\otimes1$ is an
$\mathcal{A}$-homomorphism $\Psi^*(\mathcal{G}_1)\to\Psi^*(\mathcal{G}_2)$,
which we shall denote by $\Psi^*(v)$. We have thus defined $\Psi^*$ as a
\emph{covariant functor} from the category of $\mathcal{B}$-modules to that
of $\mathcal{A}$-modules. Here this functor (unlike $\psi^*$) is no
longer exact in general, but only \emph{right exact}, since tensoring
with $\mathcal{A}$ is a right exact functor on the category of
$\psi^*(\mathcal{B})$-modules.

For every $x\in X$, one has
\[
  (\Psi^*(\mathcal{G}))_x
  =\mathcal{G}_{\psi(x)}\otimes_{\mathcal{B}_{\psi(x)}}\mathcal{A}_x,
\]
by virtue of (3.7.2). The support of $\Psi^*(\mathcal{G})$ is therefore contained in
$\psi^{-1}(\operatorname{Supp}\mathcal{G})$.
\end{env}

\begin{env}[4.3.2]
\label{0.4.3.2}
Let $(\mathcal{G}_\lambda)$ be an inductive system of $\mathcal{B}$-modules, and
let $\mathcal{G}=\varinjlim\mathcal{G}_\lambda$ be its inductive limit. The
canonical homomorphisms $\mathcal{G}_\lambda\to\mathcal{G}$ define
$\psi^*(\mathcal{B})$-homomorphisms
$\psi^*(\mathcal{G}_\lambda)\to\psi^*(\mathcal{G})$, which give a
canonical homomorphism
\[
  \varinjlim\psi^*(\mathcal{G}_\lambda)
  \longrightarrow
  \psi^*(\mathcal{G}).
\]
Since the stalk at a point of an inductive limit of sheaves is the
inductive limit of the stalks at that point (G, II, 1.11), the preceding
canonical homomorphism is \emph{bijective} (3.7.2). Moreover, the tensor product commutes
with inductive limits of sheaves, and one therefore has a \emph{canonical
functorial isomorphism}
\[
  \varinjlim\Psi^*(\mathcal{G}_\lambda)
  \xrightarrow{\sim}
  \Psi^*\bigl(\varinjlim\mathcal{G}_\lambda\bigr)
\]
of $\mathcal{A}$-modules.

On the other hand, for a finite direct sum $\bigoplus_i\mathcal{G}_i$ of
$\mathcal{B}$-modules, it is clear that
\[
  \psi^*\bigl(\bigoplus_i\mathcal{G}_i\bigr)
  =\bigoplus_i\psi^*(\mathcal{G}_i)~;
\]
hence, by tensor product with $\mathcal{A}_{[\theta]}$,
\[
  \Psi^*\bigl(\bigoplus_i\mathcal{G}_i\bigr)
  =\bigoplus_i\Psi^*(\mathcal{G}_i).
  \tag{4.3.2.1}\label{0.4.3.2.1}
\]
Passing to the inductive limit, one deduces from this, by virtue of the foregoing,
that the preceding equality still holds for an \emph{arbitrary}
direct sum.
\end{env}

\begin{env}[4.3.3]\label{0.4.3.3}
\EGAReaderTarget{canon-ref.0.4.3.3-fr}{0.4.3.3-fr}
Let $\mathcal{G}_1$, $\mathcal{G}_2$ be two $\mathcal{B}$-modules; from the
definition of inverse images of sheaves of abelian groups (3.7.1),
one immediately deduces a canonical homomorphism
\[
  \psi^*(\mathcal{G}_1)\otimes_{\psi^*(\mathcal{B})}\psi^*(\mathcal{G}_2)
  \longrightarrow
  \psi^*(\mathcal{G}_1\otimes_{\mathcal{B}}\mathcal{G}_2)
\]
of $\psi^*(\mathcal{B})$-modules, and since the stalk at a point of a
tensor product of sheaves is the tensor product of the stalks at that point
(G, II, 2.8), one deduces from (3.7.2) that the preceding homomorphism is in fact
an \emph{isomorphism}. By tensor product with $\mathcal{A}$, one deduces from it
a \emph{canonical functorial isomorphism}
\[
  \Psi^*(\mathcal{G}_1)\otimes_{\mathcal{A}}\Psi^*(\mathcal{G}_2)
  \xrightarrow{\sim}
  \Psi^*(\mathcal{G}_1\otimes_{\mathcal{B}}\mathcal{G}_2).
  \tag{4.3.3.1}\label{0.4.3.3.1}
\]
\end{env}

\begin{env}[4.3.4]
\label{0.4.3.4}
Let $\mathcal{C}$ be a $\mathcal{B}$-algebra; giving the Algebra
structure on $\mathcal{C}$ amounts to giving a $\mathcal{B}$-homomorphism
\[
  \mathcal{C}\otimes_{\mathcal{B}}\mathcal{C}\longrightarrow\mathcal{C}
\]
satisfying the conditions of associativity and commutativity (conditions
that are verified on each stalk); the preceding isomorphism makes it possible to
transport this homomorphism into a homomorphism of $\mathcal{A}$-modules
\[
  \Psi^*(\mathcal{C})\otimes_{\mathcal{A}}\Psi^*(\mathcal{C})
  \longrightarrow\Psi^*(\mathcal{C})
\]
satisfying the same conditions, so that $\Psi^*(\mathcal{C})$ is thereby endowed
with a structure of $\mathcal{A}$-algebra. In particular, it follows at once
from the definitions that the $\mathcal{A}$-algebra $\Psi^*(\mathcal{B})$ is equal to
$\mathcal{A}$ (up to a canonical isomorphism).

Likewise, if $\mathcal{M}$ is a $\mathcal{C}$-module, giving this
Module structure amounts
\oldpage[0\textsubscript{I}]{42}
to giving a $\mathcal{B}$-homomorphism
$\mathcal{C}\otimes_{\mathcal{B}}\mathcal{M}\to\mathcal{M}$ satisfying the
associativity condition; whence, by transport, a structure of
$\Psi^*(\mathcal{C})$-module on $\Psi^*(\mathcal{M})$.
\end{env}

\begin{env}[4.3.5]
\label{0.4.3.5}
Let $\mathcal{J}$ be a sheaf of ideals of $\mathcal{B}$; since the functor
$\psi^*$ is exact, the $\psi^*(\mathcal{B})$-module $\psi^*(\mathcal{J})$
is canonically identified with a sheaf of ideals of $\psi^*(\mathcal{B})$;
the canonical injection $\psi^*(\mathcal{J})\to\psi^*(\mathcal{B})$ then gives
a homomorphism of $\mathcal{A}$-modules
\[
  \Psi^*(\mathcal{J})
  =\psi^*(\mathcal{J})\otimes_{\psi^*(\mathcal{B})}\mathcal{A}_{[\theta]}
  \longrightarrow\mathcal{A}~;
\]
we shall denote by $\Psi^*(\mathcal{J})\mathcal{A}$, or by
$\mathcal{J}\mathcal{A}$ if no confusion is to be feared, the image of
$\Psi^*(\mathcal{J})$ under this homomorphism. Thus, by definition,
\[
  \mathcal{J}\mathcal{A}
  =\theta^\sharp\bigl(\psi^*(\mathcal{J})\bigr)\mathcal{A}
\]
and in particular, for every $x\in X$,
\[
  (\mathcal{J}\mathcal{A})_x
  =\theta_x^\sharp(\mathcal{J}_{\psi(x)})\mathcal{A}_x,
\]
taking into account the canonical identification of the stalks of
$\psi^*(\mathcal{J})$ with those of $\mathcal{J}$ (3.7.2). If
$\mathcal{J}_1$, $\mathcal{J}_2$ are two ideals of $\mathcal{B}$, one has
\[
  (\mathcal{J}_1\mathcal{J}_2)\mathcal{A}
  =\mathcal{J}_1(\mathcal{J}_2\mathcal{A})
  =(\mathcal{J}_1\mathcal{A})(\mathcal{J}_2\mathcal{A}).
\]
If $\mathcal{F}$ is an $\mathcal{A}$-module, we shall set
$\mathcal{J}\mathcal{F}=(\mathcal{J}\mathcal{A})\mathcal{F}$.
\end{env}

\begin{env}[4.3.6]
\label{0.4.3.6}
Let $(Z,\mathcal{C})$ be a third ringed space, and
$\Psi'=(\psi',\theta')$ a morphism $(Y,\mathcal{B})\to(Z,\mathcal{C})$;
if $\Psi''$ is the composite morphism $\Psi'\circ\Psi$, it follows from the
definition (4.3.1) and from (4.3.3.1) that
${\Psi''}^*=\Psi^*\circ{\Psi'}^*$.
\end{env}

\subsection{Relations between direct images and inverse images}
\label{subsection:0.4.4}

\begin{env}[4.4.1]
\label{0.4.4.1}
The hypotheses and notation being those of (4.2.1), let $\mathcal{G}$ be a
$\mathcal{B}$-module. By definition, a homomorphism
$u:\mathcal{G}\to\Psi_*(\mathcal{F})$ of $\mathcal{B}$-modules is
also called a \emph{$\Psi$-morphism from $\mathcal{G}$ to $\mathcal{F}$}, or
simply a \emph{homomorphism from $\mathcal{G}$ to $\mathcal{F}$}, and one
writes it $u:\mathcal{G}\to\mathcal{F}$ when no confusion results.
Giving such a homomorphism amounts to giving, for every pair $(U,V)$,
where $U$ is an open set of $X$ and $V$ an open set of $Y$ such that
$\psi(U)\subset V$, a \emph{homomorphism}
\[
  u_{U,V}:\Gamma(V,\mathcal{G})\longrightarrow\Gamma(U,\mathcal{F})
\]
\emph{of $\Gamma(V,\mathcal{B})$-modules}, $\Gamma(U,\mathcal{F})$ being
regarded as a $\Gamma(V,\mathcal{B})$-module by means of the ring
homomorphism
$\theta_{U,V}:\Gamma(V,\mathcal{B})\to\Gamma(U,\mathcal{A})$; the $u_{U,V}$
must moreover make the diagrams (3.5.1.1) commutative. Furthermore, to
define $u$ it suffices to give the $u_{U,V}$ when $U$ (resp.
$V$) ranges over a base $\mathfrak{B}$ (resp. $\mathfrak{B}'$) of the topology
of $X$ (resp. $Y$), and to verify the commutativity of (3.5.1.1) for these
restrictions.
\end{env}

\begin{env}[4.4.2]
\label{0.4.4.2}
Under the hypotheses of (4.2.1) and (4.2.6), let $\mathcal{H}$ be a
$\mathcal{C}$-module and $v:\mathcal{H}\to\Psi'_*(\mathcal{G})$ a
$\Psi'$-morphism; then
\[
  w:\mathcal{H}\xrightarrow{v}\Psi'_*(\mathcal{G})
  \xrightarrow{\Psi'_*(u)}\Psi'_*(\Psi_*(\mathcal{F}))
\]
is a $\Psi''$-morphism, which is called the \emph{composite of $u$ and $v$}.
\end{env}

\begin{env}[4.4.3]
\label{0.4.4.3}
We shall now see that one can define a canonical \emph{isomorphism}
of \emph{bifunctors} in $\mathcal{F}$ and $\mathcal{G}$
\[
  \operatorname{Hom}_{\mathcal{A}}(\Psi^*(\mathcal{G}),\mathcal{F})
  \xrightarrow{\sim}
  \operatorname{Hom}_{\mathcal{B}}(\mathcal{G},\Psi_*(\mathcal{F}))
  \tag{4.4.3.1}\label{0.4.4.3.1}
\]
which we shall denote by $v\mapsto v_\theta^\flat$ (or simply
$v\mapsto v^\flat$ if no confusion is possible); we shall denote the
inverse isomorphism by $u\mapsto u_\theta^\sharp$, or $u\mapsto u^\sharp$.
This definition is as follows: composing
$v:\Psi^*(\mathcal{G})\to\mathcal{F}$ with the canonical map
$\psi^*(\mathcal{G})\to\Psi^*(\mathcal{G})$, one obtains a homomorphism of
sheaves of groups $v':\psi^*(\mathcal{G})\to\mathcal{F}$, which is also a
homomorphism of $\psi^*(\mathcal{B})$-modules. From it one deduces (3.7.1) a
homomorphism
${v'}^\flat:\mathcal{G}\to\psi_*(\mathcal{F})=\Psi_*(\mathcal{F})$, which is
also a homomorphism of $\mathcal{B}$-modules, as one
\oldpage[0\textsubscript{I}]{43}
verifies without difficulty; one takes $v_\theta^\flat={v'}^\flat$. Likewise, from
$u:\mathcal{G}\to\Psi_*(\mathcal{F})$, which is a homomorphism of
$\mathcal{B}$-modules, one deduces (3.7.1) a homomorphism
$u^\sharp:\psi^*(\mathcal{G})\to\mathcal{F}$ of
$\psi^*(\mathcal{B})$-modules, whence, by tensoring with $\mathcal{A}$, a
homomorphism of $\mathcal{A}$-modules
$\Psi^*(\mathcal{G})\to\mathcal{F}$, which we shall denote by
$u_\theta^\sharp$. It is immediate to verify that
$(u_\theta^\sharp)_\theta^\flat=u$ and
$(v_\theta^\flat)_\theta^\sharp=v$, as well as the functoriality in
$\mathcal{F}$ of the isomorphism $v\mapsto v_\theta^\flat$. The
functoriality in $\mathcal{G}$ of $u\mapsto u_\theta^\sharp$ then follows
formally as in (3.5.4) (an argument which would also prove the
functoriality of $\Psi^*$ established directly in (4.3.1)).

If one takes for $v$ the identity homomorphism of $\Psi^*(\mathcal{G})$,
then $v_\theta^\flat$ is a homomorphism
\[
  \rho_{\mathcal{G}}:\mathcal{G}\longrightarrow
  \Psi_*(\Psi^*(\mathcal{G}))~;
  \tag{4.4.3.2}\label{0.4.4.3.2}
\]
if one takes for $u$ the identity homomorphism of $\Psi_*(\mathcal{F})$,
then $u_\theta^\sharp$ is a homomorphism
\[
  \sigma_{\mathcal{F}}:\Psi^*(\Psi_*(\mathcal{F}))
  \longrightarrow\mathcal{F}~;
  \tag{4.4.3.3}\label{0.4.4.3.3}
\]
these homomorphisms will be called \emph{canonical}. They are in general neither
injective nor surjective. There are canonical factorizations analogous to
(3.5.3.3) and (3.5.4.4).

Note that if $s$ is a section of $\mathcal{G}$ over an open set
$V$ of $Y$, $\rho_{\mathcal{G}}(s)$ is the section $s'\otimes1$ of
$\Psi^*(\mathcal{G})$ over $\psi^{-1}(V)$, where $s'$ is such that
$s'_x=s_{\psi(x)}$ for every $x\in\psi^{-1}(V)$. Note also that if
$u:\mathcal{G}\to\psi_*(\mathcal{F})$ is a homomorphism, it defines for
every $x\in X$ a homomorphism
$u_x:\mathcal{G}_{\psi(x)}\to\mathcal{F}_x$ on the stalks, obtained by
composing
$(u^\sharp)_x:(\Psi^*(\mathcal{G}))_x\to\mathcal{F}_x$ with the canonical
homomorphism $s_x\mapsto s_x\otimes1$ from $\mathcal{G}_{\psi(x)}$ to
$(\Psi^*(\mathcal{G}))_x
=\mathcal{G}_{\psi(x)}\otimes_{\mathcal{B}_{\psi(x)}}\mathcal{A}_x$.
The homomorphism $u_x$ is also obtained by passage to the inductive limit
from the homomorphisms
\[
  \Gamma(V,\mathcal{G})\xrightarrow{u}
  \Gamma(\psi^{-1}(V),\mathcal{F})\longrightarrow\mathcal{F}_x,
\]
where $V$ ranges over the neighbourhoods of $\psi(x)$.
\end{env}

\begin{env}[4.4.4]
\label{0.4.4.4}
Let $\mathcal{F}_1$, $\mathcal{F}_2$ be $\mathcal{A}$-modules,
$\mathcal{G}_1$, $\mathcal{G}_2$ $\mathcal{B}$-modules, and $u_i$ ($i=1,2$)
a homomorphism from $\mathcal{G}_i$ to $\mathcal{F}_i$. We shall denote by
$u_1\otimes u_2$ the homomorphism
\[
  u:\mathcal{G}_1\otimes_{\mathcal{B}}\mathcal{G}_2
  \longrightarrow
  \mathcal{F}_1\otimes_{\mathcal{A}}\mathcal{F}_2
\]
such that $u^\sharp=(u_1)^\sharp\otimes(u_2)^\sharp$ (taking
(4.3.3.1) into account); one verifies that $u$ is also the composite
\[
  \mathcal{G}_1\otimes_{\mathcal{B}}\mathcal{G}_2
  \longrightarrow
  \Psi_*(\mathcal{F}_1)\otimes_{\mathcal{B}}\Psi_*(\mathcal{F}_2)
  \longrightarrow
  \Psi_*(\mathcal{F}_1\otimes_{\mathcal{A}}\mathcal{F}_2),
\]
where the first arrow is the ordinary tensor product
$u_1\otimes_{\mathcal{B}}u_2$ and the second the canonical homomorphism
(4.2.2.1).
\end{env}

\begin{env}[4.4.5]
\label{0.4.4.5}
Let $(\mathcal{G}_\lambda)_{\lambda\in L}$ be an inductive system of
$\mathcal{B}$-modules, and, for every $\lambda\in L$, let $u_\lambda$ be a
homomorphism $\mathcal{G}_\lambda\to\Psi_*(\mathcal{F})$, these forming an
inductive system; set
$\mathcal{G}=\varinjlim\mathcal{G}_\lambda$ and
$u=\varinjlim u_\lambda$; then the $(u_\lambda)^\sharp$ form an
inductive system of homomorphisms $\Psi^*(\mathcal{G}_\lambda)\to\mathcal{F}$, and the
inductive limit of this system is none other than $u^\sharp$.
\end{env}

\begin{env}[4.4.6]
\label{0.4.4.6}
Let $\mathcal{M}$, $\mathcal{N}$ be two $\mathcal{B}$-modules, $V$ an open set
of $Y$, and $U=\psi^{-1}(V)$; the map $v\mapsto\Psi^*(v)$ is a
homomorphism
\[
  \operatorname{Hom}_{\mathcal{B}|V}(\mathcal{M}|V,\mathcal{N}|V)
  \longrightarrow
  \operatorname{Hom}_{\mathcal{A}|U}
  (\Psi^*(\mathcal{M})|U,\Psi^*(\mathcal{N})|U)
\]
\begin{sloppypar}
for the $\Gamma(V,\mathcal{B})$-module structures
($\operatorname{Hom}_{\mathcal{A}|U}
(\Psi^*(\mathcal{M})|U,\allowbreak\Psi^*(\mathcal{N})|U)$ is naturally endowed with a
structure of $\Gamma(U,\psi^*(\mathcal{B}))$-module, and by means of the
canonical homomorphism
\oldpage[0\textsubscript{I}]{44}
(3.7.2)
$\Gamma(V,\mathcal{B})\to\Gamma(U,\psi^*(\mathcal{B}))$, it is therefore also a
$\Gamma(V,\mathcal{B})$-module).
\end{sloppypar}
One sees at once that these homomorphisms are
compatible with restrictions, and consequently define a canonical
functorial homomorphism
\[
  \gamma:\mathcal{H}\!om_{\mathcal{B}}(\mathcal{M},\mathcal{N})
  \longrightarrow
  \Psi_*\bigl(\mathcal{H}\!om_{\mathcal{A}}
  (\Psi^*(\mathcal{M}),\Psi^*(\mathcal{N}))\bigr)~;
\]
consequently, there also corresponds to this homomorphism the homomorphism
\[
  \gamma^\sharp:
  \Psi^*\bigl(\mathcal{H}\!om_{\mathcal{B}}(\mathcal{M},\mathcal{N})\bigr)
  \longrightarrow
  \mathcal{H}\!om_{\mathcal{A}}
  (\Psi^*(\mathcal{M}),\Psi^*(\mathcal{N}))
\]
and these canonical homomorphisms are functorial in $\mathcal{M}$ and
$\mathcal{N}$.
\end{env}

\begin{env}[4.4.7]
\label{0.4.4.7}
Suppose that $\mathcal{F}$ (resp. $\mathcal{G}$) is an
$\mathcal{A}$-algebra (resp. a $\mathcal{B}$-algebra). If
$u:\mathcal{G}\to\Psi_*(\mathcal{F})$ is a homomorphism of
$\mathcal{B}$-algebras, then $u^\sharp$ is a homomorphism
$\Psi^*(\mathcal{G})\to\mathcal{F}$ of $\mathcal{A}$-algebras; this follows
from the commutativity of the diagram
\[
\begin{tikzcd}[column sep=huge,row sep=large]
\mathcal{G}\otimes_{\mathcal{B}}\mathcal{G}
  \arrow[r] \arrow[d]
  & \mathcal{G} \arrow[d,"u"] \\
\Psi_*(\mathcal{F}\otimes_{\mathcal{A}}\mathcal{F})
  \arrow[r]
  & \Psi_*(\mathcal{F})
\end{tikzcd}
\]
and from (4.4.4). Likewise, if $v:\Psi^*(\mathcal{G})\to\mathcal{F}$ is a
homomorphism of $\mathcal{A}$-algebras, then
$v^\flat:\mathcal{G}\to\Psi_*(\mathcal{F})$ is a homomorphism of
$\mathcal{B}$-algebras.
\end{env}

\begin{env}[4.4.8]
\label{0.4.4.8}
Let $(Z,\mathcal{C})$ be a third ringed space,
$\Psi'=(\psi',\theta')$ a morphism $(Y,\mathcal{B})\to(Z,\mathcal{C})$, and
$\Psi'':(X,\mathcal{A})\to(Z,\mathcal{C})$ the composite morphism
$\Psi'\circ\Psi$. Let $\mathcal{H}$ be a $\mathcal{C}$-module and $v'$ a
homomorphism from $\mathcal{H}$ to $\mathcal{G}$; the composite
$v''=v\circ v'$ is by definition the homomorphism from $\mathcal{H}$ to
$\mathcal{F}$ defined by
\[
  \mathcal{H}\xrightarrow{v'}\Psi'_*(\mathcal{G})
  \xrightarrow{\Psi'_*(v)}\Psi'_*(\Psi_*(\mathcal{F}))~;
\]
one verifies that ${v''}^\sharp$ is the homomorphism
\[
  \Psi^*({\Psi'}^*(\mathcal{H}))
  \xrightarrow{\Psi^*({v'}^\sharp)}
  \Psi^*(\mathcal{G})
  \xrightarrow{v^\sharp}\mathcal{F}.
\]
\end{env}

